\documentclass[11pt,a4paper]{article}
\usepackage[margin=2.35cm]{geometry}
\usepackage{amsmath,amssymb,booktabs,array,hyperref}
\newcommand{\Lc}{\textbf{(c)}\ }
\newcommand{\Ld}{\textbf{(d)}\ }
\newcommand{\dd}{\mathrm d}
\newcommand{\id}{\texttt{frame\_\allowbreak identities.py}}
\newcommand{\ed}{\texttt{frame\_\allowbreak edge.py}}
\newcommand{\an}{\texttt{frame\_\allowbreak analysis.py}}
\emergencystretch=3em
\title{Physical momentum, energy, and the remaining quantum frame obstruction\\in the Toda--Boussinesq corner}
\author{Quantum Boussinesq solitons project}
\date{4 November 2026}
\begin{document}
\maketitle
\begin{abstract}
\Lc\ We derive an energy counterpart of the endpoint momentum identity,
relate the endpoint responses to derivatives of dressed momentum and energy,
and give an independently implemented finite-sea dictionary and numerical
energy dispersion at fixed dressed momentum. A mixed rectangle of the
classical-frame symmetric remainder stays nonzero under the tested edge
extrapolations. These are consequences of the existing asymptotic Bethe
integral equations, not a new proof of their microscopic applicability.
\Ld\ The requested physical scalar one-loop scattering prediction is not
determined: the endpoint identities leave an antisymmetric two-body frame
function free. We exhibit an explicit freedom that survives the nested
limit while leaving every one-body observable and reciprocity test intact.
This is a checked partial result, not a completed common quantum dictionary.
No Wiener--Hopf dispersion or phase formula from Phase 2i is assumed.
\end{abstract}

\section{Scope, conventions, and input status}
Labels (a), (b), (c), (d) have the repository meanings: literature proof,
literature heuristic, project derivation/check, and guess or unresolved
identification. Numerical extrapolations labelled (c) below are explicitly
numerical evidence; they are not proofs of existence of the limiting functions.
All new formulas are checked by \id, \ed, or \an, as indicated by section.
No external source is cited or claimed to have been read. The full texts in
\texttt{literature/} were unavailable. The inputs read were the master report
(Route D and ranked priorities), \texttt{toda\_route.tex},
\texttt{toda\_edge.tex}, \texttt{matching.tex}, the Toda sections of
\texttt{charge\_frame\_master.tex}, \texttt{independent\_phase2h\_audit.tex},
\texttt{audit\_toda\_dictionary.tex}, Section 4 of
\texttt{audit\_frames\_defects.tex}, their specified implementations and saved
JSON data, and the relevant continuum formulas in \texttt{oneloop\_phase.tex}.
The overstatements in earlier notes are not adopted here.

\paragraph{Inherited conventions.}
Use $h=\hbar_T$, $t=2\pi h$, mean lattice spacing $a$, and $b=e^{-a/2}$.
The physical coordinate of the unperturbed chain is $x_T=an$.
The inherited leading classical dictionary is
\begin{equation}
 \kappa=\sqrt3\beta,\quad h=3b\hbar_B,\quad
 x_B=2\sqrt3 n,\quad t_B=2\sqrt3 b t_T,\qquad
 P_B=\frac{a\overline P}{6\sqrt3 b},\quad
 E_B=\frac{\overline E}{6\sqrt3 b^2}.
 \label{eq:units}
\end{equation}
\Lc\ The charge scalings agree with the small-amplitude soliton expansions
(\id); they are leading continuum scalings, not exact finite-amplitude
identities or a derivation of quantum collective charges.
The bad-equation input is
\[
 \begin{gathered}
 u_{tt}=u_{xx}+u_{xxxx}+3(u^2)_{xx},\quad k=\sinh\beta,\quad v=\epsilon\cosh\beta,\\
 P_{B,0}=\tfrac23\epsilon\sinh^3\beta\cosh\beta,\qquad
 E_{B,0}=\tfrac2{15}\sinh^3\beta(5+4\sinh^2\beta).
 \end{gathered}
\]
The good convention instead has $u_{tt}=u_{xx}-u_{xxxx}-3(u^2)_{xx}$,
$k=\sin\alpha$, $v=\epsilon\cos\alpha$,
$P_{g,0}=\frac23\epsilon\sin^3\alpha\cos\alpha$.
Both use $u=2(\log\tau)_{xx}$ and the stipulated symplectic normalization.
A larger right-moving bad/Toda soliton is faster; a smaller good soliton is
faster. \Lc\ This follows by differentiating the displayed velocities.
\Ld\ No real quantum equivalence between good and bad models, or analytic
continuation of their quantum states, follows from this classical dictionary.
We keep the inherited classical coordinate scale; we do not rescale time to
set the quantum sound velocity to one.

\section{An exact finite-sea one-body dictionary}
\paragraph{Definitions and derivation \Lc\ (\id, \ed).}
For an invertible finite-sea operator on $[-q,q]$, write
\begin{align}
 \Theta(d)&=2h\operatorname{Im}\log\Gamma(1+id/h)+2d\log h,
 &K(d)&=\Theta'(d),\\
 (\mathcal Af)(x)&=tf(x)-\int_{-q}^q K(x-z)f(z)\,\dd z,
 &\mathcal A\rho&=1,\quad \mathcal A\eta=x,\quad \int\rho=1/a,\\
 \mathcal AF_p(x)&=\Theta(p-x),&
 P(p)&=p+\int\Theta(p-x)\rho(x)\,\dd x,\\
 E(p)&=\frac{p^2}{2}+\int xF_p(x)\,\dd x
       =\frac{p^2}{2}+\int\Theta(p-x)\eta(x)\,\dd x .
 \label{eq:fredholm}
\end{align}
The last equality uses self-adjointness of $\mathcal A$. Thus only two sea
solutions, $\rho$ and $\eta$, are needed for all one-body observables.
For the right edge define
\begin{equation}
 \begin{gathered}
 \overline P(p)=P(p)-P(q),\qquad \overline E(p)=E(p)-E(q),\\
 D(p)=1+\int K(p-x)\rho(x)\,\dd x,\qquad
 V(p)=p+\int K(p-x)\eta(x)\,\dd x .
 \end{gathered}
 \label{eq:DV}
\end{equation}
\Lc\ Then $P'=D$, $E'=V$, $P(q)=\pi h/a$,
$D(q)=t\rho(q)$ and $V(q)=t\eta(q)$. The edge-subtracted energy uses the
existing zero-energy edge reference; no chemical-potential constant is fitted.
Parity gives $P(-p)=-P(p)$ and $E(-p)=E(p)$, so the left-edge version follows
with signed momentum. The excitation group velocity and right sound speed are
\begin{equation}
 v_{\rm exc}(p)=\frac{V(p)}{D(p)},\qquad
 v_s=\frac{\eta(q)}{\rho(q)}.
 \label{eq:velocity}
\end{equation}
These are velocities conjugate to the dressed physical momentum, not the
bare Bethe parameter $p$.

Integration by parts in the differentiated $F_p$ equation gives
\begin{equation}
 \boxed{t\rho(q)[F_p(q)-F_p(-q)]=1-D(p),\qquad
 t\eta(q)[F_p(q)+F_p(-q)]=P(p)-V(p).}
 \label{eq:newenergy}
\end{equation}
For example, multiply $\mathcal A\partial_xF_p$ by $\eta$ and use
$\int x\partial_xF_p=q[F_p(q)+F_p(-q)]-\int F_p$.
The explicit endpoint terms cancel the bare $q$ terms and leave
$E'(p)=p+\int F_p-t\eta(q)[F_p(q)+F_p(-q)]$; self-adjointness gives
$p+\int F_p=P(p)$. This proves the second identity without assuming a
Wiener--Hopf factorization. \Lc

\paragraph{Exact quantum Jacobian \Lc\ (\id, \ed).}
The old variable $y$ is best understood as a monotone \emph{physical-momentum
label}, not a bare spectral coordinate:
\begin{equation}
 \overline P(q+hu(y))=P_0(\kappa)
 =\frac{4b}{a}(\kappa\cosh\kappa-\sinh\kappa),\qquad
 \kappa=\sqrt{yh/b}.
 \label{eq:invert}
\end{equation}
Consequently the exact finite-$h$ dictionary is the parametric pair
$(\overline P(q+hu),\overline E(q+hu))$, with
\begin{equation}
 \boxed{\frac{\dd u}{\dd y}=\frac{2\sinh\kappa}{aD(q+hu)},\qquad
 \frac{\dd\overline E}{\dd y}
 =\frac{2h\sinh\kappa}{a}\frac{V(q+hu)}{D(q+hu)}.}
 \label{eq:jac}
\end{equation}
The classical comparator is $u_{0,h}(y)=2b(\cosh\kappa-1)/h$ and
$u'_{0,h}=\sinh\kappa/\kappa$. Subtracting these known finite-amplitude
expressions before extrapolating removes a large classical lattice artifact.
All new calculations use the exact endpoint $q$, with no $q(1+10^{-10})$
offset. Fixed nonzero $\overline P$ as $h\to0$ and fixed $y$ are different
limits: at fixed $y$, $a\sqrt b\,\overline P/h^{3/2}\to\frac43 y^{3/2}$.
The term ``fixed physical momenta in the edge limit'' below means fixed
\emph{scaled} dressed momenta specified by \eqref{eq:invert}.

\section{Energy dispersion: what can be extracted without a phase frame}
\paragraph{Classical comparison and loop order \Lc\ (\id).}
The classical energy is $E_0(\kappa)=2b^2(\sinh\kappa\cosh\kappa-\kappa)$;
$dE_0/dP_0=ab\sinh\kappa/\kappa$. Moreover
\begin{equation}
 E_0-abP_0=\frac{2b^2}{15}\kappa^5+O(\kappa^7),\qquad
 E_{B,0}-P_{B,0}=\frac15\beta^5+O(\beta^7).
 \label{eq:dispcl}
\end{equation}
At fixed spectral coordinate, suppose locally away from the edge that
$\overline P=P_0(p)+hP_1(p)+o(h)$ and
$\overline E=E_0(p)+hE_1(p)+o(h)$. Inverting at fixed $P=P_0(p_0)$ gives
\begin{equation}
 p=p_0-h\frac{P_1(p_0)}{P'_0(p_0)}+o(h),\qquad
 E(P)=E_0(p_0)+h\left[E_1(p_0)-\frac{E'_0(p_0)}{P'_0(p_0)}P_1(p_0)\right]+o(h).
 \label{eq:fixedP}
\end{equation}
This is an exact formal chain-rule statement conditional on the expansions,
not a uniform endpoint expansion. An $O(h)$ energy correction is the
one-loop correction to the energy, whereas the one-loop \emph{phase} is
$O(h^0)$ after division of an action by $h$.

\paragraph{Independent numerical data \Lc\ (\ed, \an).}
At $a=1$ the calculation supplies a useful finite-$y$ coefficient
\begin{equation}
 c_{E,h}(y)=\frac{\overline E(P_0(\sqrt{yh/b}))-E_0(\sqrt{yh/b})}{hP_0(\sqrt{yh/b})}.
 \label{eq:ce}
\end{equation}
A finite extrapolated $c_E(y)$ would give
$\overline E=E_0(P)+hP c_E(y)+o(hP)$ at fixed $y$.
It contains the quantum crossover at that $y$; it is not itself a
large-$y$ loop expansion. Table~\ref{tab:one} records the independent fits.
In the inherited units \eqref{eq:units} its leading quantum contribution to
bad-Boussinesq energy is $3\hbar_B P_B c_E(y)$ for $a=1$.
\Ld\ This is a Toda-induced small-amplitude dispersion, not a comparison
with a fully renormalized continuum energy operator. Quantum corrections to
the identification of collective charges remain a separate question.


\par
\begin{table}[ht]\centering\small
\setlength{\tabcolsep}{4pt}
\begin{tabular}{rrrrrrr}\toprule
$y$ & $u-y$ & spread & $u\prime$ & spread & $c_E(y)$ & spread\\\midrule
0.1 & -0.0763583 & $6.9\,10^{-9}$ & 0.3524490 & $9.6\,10^{-8}$ & 0.0725266 & $5.9\,10^{-9}$\\
0.5 & -0.2505729 & $5.6\,10^{-8}$ & 0.7041319 & $1.3\,10^{-7}$ & 0.0601062 & $2.3\,10^{-8}$\\
1 & -0.3574864 & $1.3\,10^{-7}$ & 0.8429770 & $1.8\,10^{-7}$ & 0.0559987 & $7.3\,10^{-8}$\\
2 & -0.4663924 & $3.7\,10^{-7}$ & 0.9213452 & $2.8\,10^{-7}$ & 0.0543111 & $2.2\,10^{-7}$\\
4 & -0.5757147 & $1.1\,10^{-6}$ & 0.9604744 & $4.4\,10^{-7}$ & 0.0536443 & $7.2\,10^{-7}$\\
8 & -0.6854957 & $3.5\,10^{-6}$ & 0.9801739 & $7.5\,10^{-7}$ & 0.0533400 & $2.5\,10^{-6}$\\
16 & -0.7955379 & $1.2\,10^{-5}$ & 0.9900702 & $1.3\,10^{-6}$ & 0.0531937 & $8.8\,10^{-6}$\\
32 & -0.9057144 & $4.2\,10^{-5}$ & 0.9950310 & $2.4\,10^{-6}$ & 0.0531202 & $3.2\,10^{-5}$\\
64 & -1.0159485 & $1.5\,10^{-4}$ & 0.9975150 & $4.2\,10^{-6}$ & 0.0530723 & $1.1\,10^{-4}$\\
\bottomrule\end{tabular}
\caption{Numerical edge extrapolations at fixed scaled dressed momentum (c). The spreads concern edge-fit models only. Maximum measured mesh differences across these rows are $5.1\,10^{-10}$ for $u-y$, $1.9\,10^{-11}$ for $u\prime$, and $4.7\,10^{-10}$ for $c_E$. Their propagated mesh diagnostics are at most fifteen times these values. No large-$y$ asymptotic coefficient is fitted here.}\label{tab:one}
\end{table}
\par


\Lc\ At fixed, unscaled momenta the independent finite-$h$ extrapolation
also reproduces the earlier dispersion estimates; Table~\ref{tab:fixed}
reports the coefficient $[\overline E(P)-E_0(P)]/h$.
The sound-velocity estimate is tabulated separately and is not silently
removed by a quantum change of coordinates.

\par
\begin{table}[ht]\centering\small
\begin{tabular}{rrrr}\toprule
$P$ & $E_1(P)$ & edge-fit spread & mesh difference\\\midrule
0.15 & 0.007775546 & $7.5\,10^{-8}$ & $8.9\,10^{-12}$\\
0.4 & 0.020319171 & $3.7\,10^{-8}$ & $3.9\,10^{-11}$\\
0.6 & 0.030098075 & $2.2\,10^{-8}$ & $7.1\,10^{-11}$\\
\bottomrule\end{tabular}
\caption{One-loop energy at fixed, unscaled physical momentum (c, numerical). These are finite-amplitude Toda results, not finite-amplitude Boussinesq predictions.}\label{tab:fixed}
\end{table}
\par

\Lc\ The extrapolated $(v_s-b)/h$ is $0.0795774724$, with edge-fit spread $5.7\,10^{-9}$ and measured mesh difference $1.2\,10^{-11}$. This is a numerical estimate, not an independently proved closed form.

\section{Endpoint responses: derivatives, not yet soliton charges}
\paragraph{Response dictionary \Lc\ (\id, \ed).}
Put $f_p^\pm=F_p(\pm q)-F_q(\pm q)$ and use the earlier audit's responses
\begin{equation}
 \mathfrak M_p=ta\rho(q)(f_p^+-f_p^-),\qquad
 \mathfrak Q_p=\frac{f_p^++f_p^-}{2a\rho(q)}.
\end{equation}
The new energy identity turns these into directly measurable derivatives
of the one-body Bethe dispersion:
\begin{equation}
 \boxed{\mathfrak M_p=-a[D(p)-D(q)],\qquad
 \mathfrak Q_p=\frac{\overline P(p)-V(p)+V(q)}{2at\rho(q)\eta(q)}.}
 \label{eq:responses}
\end{equation}
There is also a thermodynamic interpretation of the \emph{unsubtracted}
sum. At fixed $h$ vary $q$ and let $n=1/a=\int\rho$. Differentiating the sea
problem and using self-adjointness yields
\begin{equation}
 \partial_q n=2t\rho(q)^2,\qquad
 \partial_qP(p)=t\rho(q)[F_p(q)+F_p(-q)],\qquad
 -a\partial_a P(p)\big|_{p,h}=\frac{F_p(q)+F_p(-q)}{2a\rho(q)}.
 \label{eq:pressure}
\end{equation}
Thus $\mathfrak Q$ is a difference of these fixed-bare-parameter pressure
responses, with the second derivative evaluated at $p=q$ \emph{after}
differentiation. It is not a derivative keeping $p-q$ fixed or following a
fixed physical-momentum excitation. Those operations must not be conflated.
The checks use independent finite differences in $q$ as well as endpoint
identities, at 60 decimal digits.

\Ld\ Equations \eqref{eq:responses}--\eqref{eq:pressure} do not identify
$\mathfrak M$ with a quantum integrated compression or $\mathfrak Q$ with
$P-\Pi$ in a chosen collective-coordinate Hilbert space. Their classical
analogy alone is insufficient for that step.

\section{Two-body references and the surviving mixed diagnostic}
\paragraph{Exact decomposition \Lc\ (\id, \ed).}
Keep the directional Bethe kernel
\begin{equation}
 \Phi(p,z)=\Theta(p-z)-\int K(p-x)F_z(x)\,\dd x .
\end{equation}
Its four-term right-edge reference is
$\widehat\Phi(p,z)=\Phi(p,z)-\Phi(p,q)-\Phi(q,z)+\Phi(q,q)$.
The old two-term reference is
$N(p,z)=\Phi(p,z)-\Phi(p,q)=\widehat\Phi(p,z)+U(z)$ with
$U(z)=\Phi(q,z)-\Phi(q,q)$. Thus switching these references adds a one-body
function. The exact endpoint reciprocity identity becomes
\begin{equation}
 \widehat\Phi(p,z)+\widehat\Phi(z,p)
 =t(f_p^+f_z^+-f_p^-f_z^-)
 =\mathfrak M_p\mathfrak Q_z+\mathfrak M_z\mathfrak Q_p.
 \label{eq:recip}
\end{equation}
Our independent implementation checks it against the directly solved phases,
not only against another expression in the same response functions.
One algebraically antisymmetric representative is
\begin{equation}
 B(p,z)=\widehat\Phi(p,z)-\mathfrak M_p\mathfrak Q_z.
 \label{eq:B}
\end{equation}
\Ld\ We do not select $B/h$ as the physical scalar scattering phase.
Antisymmetry is necessary for such a proposal, but is not its derivation.

\paragraph{A reference-free numerical test \Lc\ (\ed, \an).}
At the fixed dressed momenta of \eqref{eq:invert} define
\begin{align}
 C_{\rm cl}(y,z)&=-2\kappa_y[P_0(\kappa_z)-2b\sinh\kappa_z],\\
 S_h(y,z)&=\frac{\mathfrak M_y\mathfrak Q_z+
 \mathfrak M_z\mathfrak Q_y-C_{\rm cl}(y,z)-C_{\rm cl}(z,y)}h,\\
 \mathcal D_h(Y)&=S_h(2Y,Y/2)-S_h(Y,Y/2)
                    -S_h(2Y,Y/4)+S_h(Y,Y/4).
 \label{eq:rectangle}
\end{align}
All calculations here use $a=1$ and the exact quantum $q$.
Every additive one-body reference disappears from $\mathcal D_h$.
Its values therefore test genuine mixed dependence in the retained
classical frame, and the same rectangle is obtained from the two-term
normalization. Table~\ref{tab:rect} gives the new extrapolation.

\par
\begin{table}[ht]\centering\small
\begin{tabular}{rrrr}\toprule
$Y$ & $\mathcal D_0(Y)$ estimate & edge-fit spread & mesh difference\\\midrule
4 & -0.5043623 & $3.7\,10^{-4}$ & $6.0\,10^{-11}$\\
8 & -0.4349827 & $1.3\,10^{-3}$ & $1.4\,10^{-10}$\\
16 & -0.4087054 & $5.2\,10^{-3}$ & $2.4\,10^{-10}$\\
32 & -0.3957520 & $2.2\,10^{-2}$ & $3.1\,10^{-10}$\\
\bottomrule\end{tabular}
\caption{Mixed rectangles of the classical-frame symmetric discrepancy (c, numerical). The two largest momenta are in ratio $2^{3/2}$ at leading edge order; all four corners use separately inverted dressed momenta.}\label{tab:rect}
\end{table}
\par

\Lc\ Fitting these four edge estimates to $A+B/\sqrt Y$, $A+B/Y$, or $A+B/\sqrt Y+C/Y$ gives $A=-0.3276,-0.3774,-0.4034$, respectively. The fit-model range is about $0.0758$, separately from the tabulated edge-fit errors. These unconstrained fits suggest a nonzero mixed asymptote but do not establish it; no fit is selected as the physical one-loop prediction.

\Lc\ These data support survival of the nonseparable discrepancy at fixed
scaled physical momenta as $h\to0$, at the tested finite $Y$. This is stronger
than the earlier single finite-$h$ counterexample. \Ld\ They do not prove a
nonzero $Y\to\infty$ limit. Different large-$Y$ fit models are reported in
the JSON; extrapolation of four moderately large values is not an asymptotic
theorem. Nor does a symmetric directional-kernel discrepancy disprove an
ordered forward amplitude: it constrains a proposed common scalar frame.

\paragraph{Endpoint subtraction cannot be dropped.}
\Lc\ In \eqref{eq:responses} the difference $D(p)-D(q)$ is essential.
Replacing it by $D(p)$ does not compute the four-term phase. In particular,
a formula for its symmetric part using only $u'(y)$ at the two external
points, with no subtracted endpoint response, is not a consequence of
\eqref{eq:recip}. The direct phase test in \ed\ retains both endpoints.
\Ld\ No claimed Phase 2i symmetric-limit formula is used to extrapolate
these data.

\section{Why there is still no unique one-loop phase prediction}
\paragraph{An explicit remaining freedom \Lc\ (\id, \an).}
For $y_1,y_2>0$ and any real $\lambda$, define
\begin{equation}
 A_\lambda(y_1,y_2)=\lambda\frac{y_1y_2(y_1-y_2)}{(y_1+y_2)^3}.
 \label{eq:ambiguity}
\end{equation}
It is antisymmetric, vanishes if either momentum label goes to zero with
the other fixed, and has a nonzero mixed derivative. Replacing a candidate
frame representative $B$ by $B+hA_\lambda$ changes no one-body momentum,
energy, Jacobian, classical action limit, or symmetric reciprocity identity.
Nevertheless it changes the phase at order $h^0$, with
\begin{equation}
 A_\lambda(Y,r^2Y)=\lambda\frac{r^2(1-r^2)}{(1+r^2)^3},\qquad 0<r<1,
 \label{eq:ambiguitylimit}
\end{equation}
independently of $Y$. A mixed rectangle of this term is nonzero, so this
freedom is not a one-body phase convention in disguise.
This is a counterexample to \emph{uniqueness from the one-body and reciprocity
constraints}; it is not a proposal to modify the microscopic Toda theory.

\Ld\ The missing input is a microscopic derivation of the quantum canonical
position/frame map (or equivalent two-particle quantization/delay condition)
that fixes the antisymmetric part of the transformation to the continuum
Darboux frame. It must identify the quantum compression and momentum charges
including endpoint variations. Exact one-body dispersion, or imposing scalar
antisymmetry, does not supply it. Deriving that input, together with a uniform
large-$Y$ edge analysis, is required before a continuum one-loop calculation
can be assigned a unique mandatory Toda target.

Consequently no number or table in this note is advertised as the requested
physical two-soliton one-loop phase. \Ld\ A continuum calculation using the
same inherited coordinates can compare to the one-body dispersion and must
respect the exact momentum Jacobian; its full phase cannot yet be tested
against a uniquely fixed scalar target. This is the precise obstruction in
this partial result.

\section{Limits and the conditional continuum comparison}
\paragraph{What commutes and what is not established.}
\Lc\ At finite $h$, within a smooth invertible sea problem with $D>0$, inversion
from $p$ to $P$ is a regular coordinate change: the chain rules above and
endpoint referencing can be performed exactly in either coordinate.
In the displayed ambiguity example, $h\to0$ and $Y\to\infty$ commute at
fixed $r$, since \eqref{eq:ambiguitylimit} is independent of both.
That observation is \emph{only} about the counterexample.
\Ld\ Commutation of $h\to0$ and $Y\to\infty$ for the full physical phase
has not been established. At fixed $h$, $Y\to\infty$ exits the small-amplitude
corner, since $\kappa^2=hY/b$; a simultaneous small-amplitude limit needs
$hY/b\to0$. Taking $Y\to\infty$ first is not the same continuum procedure.
\Lc\ Likewise $r\to0$ at fixed $Y$ takes the smaller excitation to its edge,
whereas a semiclassical expansion of both excitations requires $r^2Y\gg1$.
Equal amplitudes require separate analysis: the classical logarithm
$\log[(1-r)/(1+r)]$ is singular and the spectral separation need not be large
when $(1-r^2)Y$ is small. These facts obstruct a uniform termwise expansion;
they are not proofs that the exact physical amplitudes have unequal iterated
limits. \Ld\ The actual phase limit-commutation questions at $r=0,1$ remain
open in this work, including the head-on channel.

\paragraph{Comparison with \texttt{oneloop\_phase.tex}.}
\Lc\ Conditional on that note's mode-sum momentum shift,
$\Delta P_g=\epsilon\alpha(\log\alpha+c)/\pi+O(\alpha^3\log\alpha)$,
$c=-\log2-1+\frac14\log3$, the leading relative shift is
\begin{equation}
 \frac{\hbar_B\Delta P_g}{P_{g,0}}
 =\frac{3}{2\pi y}(\log\alpha+c)+\cdots,\qquad y=\alpha^2/\hbar_B .
 \label{eq:nonuniform}
\end{equation}
The algebra is independently checked in \id. The expansion is nonuniform
in the fixed-$y$ edge limit; it requires $y\gg|\log\alpha|$.
The master report also makes the physical one-loop determinant, endpoint
domain, and normalization conditional. \Ld\ Neither reversing the overtaking
ordering nor replacing that momentum shift by the proposed collective one
proves a common quantum dictionary. We therefore do not claim agreement or
inequivalence with the continuum phase integrals. No new regularization or
absolute phase convention is selected.

\section{Verification, error budget, and failures recorded}
\paragraph{Numerical protocol \Lc.}
\ed\ is a fresh finite-sea Nystr\"om implementation using SciPy digamma and
log-gamma functions, independent of the old solver classes. It solves the
count equation at fixed $a=1$ and uses LU reuse for backflow. The base meshes
are $N=325,450,700,1200,2200$ at $h=.08,.04,.02,.01,.005$;
refinements have $N=1620$ at $.01$ and $2750$ at $.005$.
The classical momentum target is inverted before every energy or phase
measurement. No result is obtained by fitting to a proposed target function.
\id\ independently uses mpmath at 60 digits, $N=40,64$, $h=.4$, $q=1.5$,
and determines $a$ from counting. It tests energy and momentum derivatives,
endpoint identities, parity samples, and pressure responses with symmetric
$q$ differences of step $10^{-10}$. Maximum tested errors are approximately
$1.31\times10^{-9}$ and $4.62\times10^{-15}$, respectively; quadrature and
finite-difference errors, not decimal precision, limit these tests.

The central edge fit is quadratic through the last three $h$ values.
The extrapolation spread is the maximum difference from a last-two linear
fit, a last-four cubic fit, and a last-four fit in $1,h,h^2,h\log h$.
The mesh column is the largest observed refined-minus-base difference at
$.01$ or $.005$. The JSON also multiplies this by the largest $\ell^1$
norm of the extrapolation weights as a conservative diagnostic of amplification.
None of these is a rigorous confidence interval or bound on omitted terms.
Large-$Y$ fit-model spread is separate from edge-extrapolation spread.
There is no assigned numerical error bar for the unknown physical frame.

\paragraph{Phase 2i use policy.}
\Lc\ The only checked Phase 2i formula used for a diagnostic is the product
of its printed factors. Direct 60-digit Gamma-function evaluation at
$|\omega|=.2,1,3$ gives
\begin{equation}
 \Phi_{\rm WH}(i\omega)\Phi_{\rm WH}(-i\omega)
 =\frac{\pi e^{|\omega|/2}}{\sinh(|\omega|/2)}
\end{equation}
to better than $10^{-55}$, and disagrees with its printed Fourier symbol.
This verifies a factor product only; zero-frequency distributional issues,
the half-line inverse, $c_{\rm edge},d_0,e_0$, Krein convolution, and $c_1(r)$
are not adopted. The new numerical dispersion and remainder data do not use
any of those formulas.

\begin{center}\small
\begin{tabular}{p{.48\linewidth}cp{.35\linewidth}}
\toprule Claim & Label & Check or limitation\\\midrule
Exact energy/backflow identity and velocity & (c) & \id; independent \ed\\
Physical-momentum inversion and quantum Jacobian & (c) & symbolic chain rule; direct solve\\
Pressure interpretation of endpoint responses & (c) & 60-digit finite differences\\
One-body fixed-$P$ energy and finite-$y$ tables & (c) & mesh and extrapolation diagnostics\\
Nonzero mixed edge remainder at tested $Y$ & (c) & direct reciprocity and rectangles\\
Remaining antisymmetric freedom & (c) & symbolic counterexample; mixed rectangle\\
Unique scalar one-loop phase; full commuting limits & (d) & not established\\
\bottomrule
\end{tabular}
\end{center}

\paragraph{Failures recorded.}
\begin{enumerate}
\item The hoped-for disappearance of the classical-frame mixed remainder
under the physical-momentum inversion is not supported by the new edge
scan. This negative result is retained, not removed by redefining a reference.
\item The printed Phase 2i Fourier symbol fails the independent factor-product
check. No other Phase 2i result is promoted to an input on that basis.
\item The attempted determination of a unique scalar one-loop target fails:
\eqref{eq:ambiguity} preserves all imposed one-body and reciprocity constraints
but changes that target. This is an obstruction of the stated derivation,
not a proof that a microscopic answer does not exist.
\item Large-$Y$ extrapolation and the endpoint limits are unresolved; the
numerical error diagnostics must not be presented as a proof of a constant
asymptote or of commutation of exact limits.
\item The first LaTeX build failed because its temporary output directory had not yet been created (no log file could be opened). Creating that directory fixed the build setup; this was not a sandbox rejection. The first successful layout also had two overlong displayed equations; both were split before the final rendering. One render was started while a rebuild was still writing the PDF and failed with an empty-stream error; it was repeated after the build completed.
\item No numerical assertion tolerance was weakened. No existing project
script or document was modified. The full historical suite was not rerun:
these deliverables contain independent checks, and no prior implementation
is changed or newly certified.
\end{enumerate}

\paragraph{Reproduction.}
Run with one BLAS thread from \texttt{phase2/scripts/}:
\begin{verbatim}
OPENBLAS_NUM_THREADS=1 OMP_NUM_THREADS=1 python3 frame_edge.py --recompute
python3 frame_identities.py
python3 frame_analysis.py
\end{verbatim}
Without \texttt{--recompute}, the edge script validates its saved JSON.
Every completed checker ends with \texttt{ALL CHECKS PASSED}; this certifies
its explicitly stated identities and diagnostics, not the unresolved physical
identifications. SHA-256 provenance and raw timings are saved in the JSON
outputs. Build from \texttt{phase2/} with \texttt{pdflatex} twice.
The recorded full edge scan takes about 326 seconds, the high-precision identity run about 178 seconds, and the cached analysis under one second on this machine. A complete regeneration is therefore about eight and a half minutes; cached edge validation avoids the dense solves.
\end{document}
